% !TEX root = ../main.tex
% !TEX program = lualatex
\documentclass[../main.tex]{subfiles}
\IfSubfilesClassLoaded{\externaldocument{\subfix{../build/main}}}{}
% =============================================================================
% CHAPTER 7: FIELD ACTIVITY FINANCIAL MANAGEMENT
% DESCRIPTION: Cost categories, General Fund vs NWCF operations, and funding documents.
% =============================================================================

\begin{document}
\IfSubfilesClassLoaded{\chapter{EDO Study Guide}}{}
%====================
% Contents here
%====================
% === FIELD ACTIVITY FINANCIAL MANAGEMENT ===
\section{Field Activity Financial Management}

\subsection{Cost categories}
\begin{description}
  \item[\textbf{Direct costs.}] Costs directly attributable to a project (for example, materials, labor, and travel tied to a specific job)~\autocite{edo-3-1-5-field-activity-fm-2025}.
  \item[\textbf{Indirect costs.}] Costs that cannot be attributed to a single project (for example, security, executive salaries, and shared services)~\autocite{edo-3-1-5-field-activity-fm-2025}.
\end{description}

\note{Direct costs typically include labor, material, and \ac{odc}; indirect costs include overhead and \ac{ganda}~\autocite{edo-3-1-5-field-activity-fm-2025}.}

\subsection{Financial management systems}
\begin{description}
  \item[\textbf{Mission funded (General Fund).}] Appropriated funds flow down the chain of command to execute mission needs; the goal is to budget adequate resources without excess~\autocite{edo-3-1-5-field-activity-fm-2025}.
  \item[\textbf{Working capital funded (NWCF).}] A reimbursable system where activities are not directly appropriated; they recover costs from customers and aim to break even over time~\autocite{edo-3-1-5-field-activity-fm-2025}.
\end{description}

\note{General Fund activity management is a \textit{``checking account''} model; \ac{nwcf} is a \textit{``credit card''} model. Some \ac{nwcf} commands have mission-funded units, but funds cannot be mixed~\autocite{edo-3-1-5-field-activity-fm-2025}.}

\subsection{Activity examples}
\begin{longtblr}[
    label   = {tab:activity_examples},
    caption = {Mission-funded vs NWCF activity examples},
    entry   = {Activity Examples},
    remark{Source} = {\tabCite{edo-3-1-5-field-activity-fm-2025}}
  ]
  {colspec = {@{} X X @{}}}
  \toprule
  {Mission funded (General Fund)} & {NWCF activities} \\
  \midrule
  \textbullet\ Naval Shipyards, Regional Maintenance Centers, Ship Repair Facilities \newline
  \textbullet\ Supervisors of Shipbuilding, Program Executive Offices \newline
  \textbullet\ Systems commands (HQ only) \newline
  \textbullet\ Operating forces (fleet commanders, TYCOMs, ships) \newline
  \textbullet\ Most other activities (e.g., EDO School)
  &
  \textbullet\ SYSCOM Echelon III activities (NAVSUP FLCs; NAVAIR/NAVSEA/NAVWAR warfare centers; NAVFAC facility engineering commands and expeditionary warfare centers; Naval Research Laboratory) \newline
  \textbullet\ Military Sealift Command \newline
  \textbullet\ Depot maintenance (non-ship): Fleet Readiness Centers (aircraft), Marine Corps depots \\
  \bottomrule
\end{longtblr}

\subsection{General Fund operations}
\begin{description}
  \item[\textbf{Appropriation flow.}] \ac{asnfmc} (Echelon I) allocates funds to Echelon II \acp{bso}; BSOs distribute allotments and \acp{eob}. Echelon III field activities issue allowances and \acp{optar} to subordinate units at Echelon IV~\autocite{edo-3-1-5-field-activity-fm-2025}.
\end{description}

\subsection{Appropriated vs reimbursable funds for mission-funded activities}
\begin{description}
  \item[\textbf{Appropriated funds.}] Pay for mission labor (civilian salaries), mission non-labor (materials, travel, contracts), and all indirect costs (overhead)~\autocite{edo-3-1-5-field-activity-fm-2025}.
  \item[\textbf{Reimbursable funds.}] Used when a mission-funded activity performs work for a customer; reimbursements cover direct labor and other direct costs, but not indirect costs~\autocite{edo-3-1-5-field-activity-fm-2025}.
\end{description}

\subsection{General Fund budgeting cycle (EOB)}
\begin{enumerate}
  \item Mid-year review of the current year budget~\autocite{edo-3-1-5-field-activity-fm-2025}.
  \item Answer budget data calls~\autocite{edo-3-1-5-field-activity-fm-2025}.
  \item Submit next year's (future) budget~\autocite{edo-3-1-5-field-activity-fm-2025}.
  \item Marks and reclamas (summer review)~\autocite{edo-3-1-5-field-activity-fm-2025}.
  \item Sept/Oct: end of current FY, start of next FY; receive \ac{eob} and execute~\autocite{edo-3-1-5-field-activity-fm-2025}.
\end{enumerate}

\subsection{What is the NWCF?}
\begin{description}
  \item[\textbf{Definition.}] \ac{nwcf} is a reimbursable operation that provides goods or services to customers and recovers the full cost of output~\autocite{edo-3-1-5-field-activity-fm-2025}.
  \item[\textbf{Criteria to join.}] Output costs must be identifiable, customers must be known, and the advantages and disadvantages of the buyer-seller relationship must be evaluated~\autocite{edo-3-1-5-field-activity-fm-2025}.
\end{description}

\subsection{NWCF cycle of operations}
\begin{enumerate}
  \item Congress provides corpus to begin operations~\autocite{edo-3-1-5-field-activity-fm-2025}.
  \item The \ac{nwcf} activity receives a reimbursable work request via a funding document~\autocite{edo-3-1-5-field-activity-fm-2025}.
  \item The activity performs work; direct and indirect costs are financed by the corpus~\autocite{edo-3-1-5-field-activity-fm-2025}.
  \item Costs are priced using stabilized rates held constant during execution~\autocite{edo-3-1-5-field-activity-fm-2025}.
  \item The customer is billed for costs incurred; payment reimburses the corpus~\autocite{edo-3-1-5-field-activity-fm-2025}.
\end{enumerate}

\subsection{Stabilized labor rates (SLR)}
\begin{description}
  \item[\textbf{Purpose.}] \ac{slr} reimburses \ac{nwcf} activities for work performed and is established during the budget process to recover full costs, prior period gains/losses, and approved capital surcharges~\autocite{edo-3-1-5-field-activity-fm-2025}.
  \item[\textbf{Formula.}]
        \[
          \text{SLR} = \frac{\text{Direct labor} + \text{Indirect (overhead/G\&A)} + \text{AOR recoupment}}{\text{Direct hours}}
        \]
  \item[\textbf{Components.}] Direct labor, indirect expenses (including \ac{ganda}), and \ac{aor} recoupment to achieve break-even~\autocite{edo-3-1-5-field-activity-fm-2025}.
\end{description}

\subsection{Operating results: NOR and AOR}
\begin{description}
  \item[\textbf{\ac{nor}.}] Current-year profit or loss based on revenue minus expenses~\autocite{edo-3-1-5-field-activity-fm-2025}.
  \item[\textbf{\ac{aor}.}] Cumulative gains and losses since inception; \ac{aor} is managed at Navy level and used to set rates~\autocite{edo-3-1-5-field-activity-fm-2025}.
  \item[\textbf{Rate adjustment.}] Gains reduce rates two years out; losses increase rates two years out to drive \ac{aor} toward zero in the budget year~\autocite{edo-3-1-5-field-activity-fm-2025}.
\end{description}

\subsection{NWCF budget timeline}
\begin{description}
  \item[\textbf{Feb--Apr.}] Internal budget formation at the activity level~\autocite{edo-3-1-5-field-activity-fm-2025}.
  \item[\textbf{Apr--May.}] Activity budgets consolidate to activity group and submit to the \ac{bso}~\autocite{edo-3-1-5-field-activity-fm-2025}.
  \item[\textbf{May--Jun.}] \ac{bso} submits the activity group budget to \ac{fmb}~\autocite{edo-3-1-5-field-activity-fm-2025}.
  \item[\textbf{Jul--Aug.}] \ac{fmb} submits the activity group budget to \ac{dod} for the \ac{bes}~\autocite{edo-3-1-5-field-activity-fm-2025}.
  \item[\textbf{Sep.}] \ac{dod} submits to Congress~\autocite{edo-3-1-5-field-activity-fm-2025}.
  \item[\textbf{Feb.}] President signs the budget (first Monday in February); rates are set 18--24 months in advance~\autocite{edo-3-1-5-field-activity-fm-2025}.
\end{description}

\subsection{Funding documents and cost stages}
\begin{description}
  \item[\textbf{Reimbursable vs direct cite.}] Reimbursable work charges costs to the \ac{nwcf} corpus and then bills the customer; direct cite (commercial contracts only) charges the customer's line of accounting directly~\autocite{edo-3-1-5-field-activity-fm-2025}.
  \item[\textbf{Cost stages.}] Committed (funding document received), obligated (contract awarded), costed (work performed/material received)~\autocite{edo-3-1-5-field-activity-fm-2025}.
\end{description}

\subsection{Carryover}
\begin{description}
  \item[\textbf{Definition.}] Dollar value of reimbursable work ordered and funded but not yet completed at FY end~\autocite{edo-3-1-5-field-activity-fm-2025}.
  \item[\textbf{Formula.}] Carryover = Available balance (cash) + commitments + obligations (costed transactions excluded)~\autocite{edo-3-1-5-field-activity-fm-2025}.
  \item[\textbf{Rules.}] Some carryover is required to start the new FY; excess carryover can trigger congressional marks and future budget reductions~\autocite{edo-3-1-5-field-activity-fm-2025}.
\end{description}

\subsection{Project orders (PO)}
\begin{description}
  \item[\textbf{Specificity.}] Description of work and period of performance must be specific and produce a tangible end product~\autocite{edo-3-1-5-field-activity-fm-2025}.
  \item[\textbf{Timing.}] Work may extend past fund expiration but not beyond cancellation; \acp{po} cannot be issued solely to extend funds and scope cannot expand after expiration~\autocite{edo-3-1-5-field-activity-fm-2025}.
  \item[\textbf{51/49 rule.}] At least 51\% of reimbursable work must be performed in-house; no more than 49\% can be out-house~\autocite{edo-3-1-5-field-activity-fm-2025}.
  \item[\textbf{Restrictions.}] \acp{po} are reimbursable only and limited to \ac{dod} activities; they extend the cited funds beyond expiration and carry additional restrictions~\autocite{edo-3-1-5-field-activity-fm-2025}.
\end{description}

\subsection{Other funding documents}
\begin{description}
  \item[\textbf{\ac{mipr}.}] Transfers funds between military organizations for supplies or services; can be accepted on a direct cite or reimbursable basis~\autocite{edo-3-1-5-field-activity-fm-2025}.
  \item[\textbf{\ac{ipr}.}] Used when requiring and performing activities are from different departments (outside \ac{dod}); cannot cite the Project Order statute but may cite Economy Act or other statutory authorities~\autocite{edo-3-1-5-field-activity-fm-2025}.
  \item[\textbf{FS Form 7600B.}] Treasury G-Invoicing form for reimbursable \acp{igt}; applicable to manual or electronic submissions and not limited to Economy Act or Project Order authorities~\autocite{edo-3-1-5-field-activity-fm-2025}.
\end{description}

\subsection{Summary check}
\begin{description}
  \item[\textbf{NWCF success.}] Financial success means revenue and expenses break even over time (\ac{aor} near zero) and pricing remains affordable~\autocite{edo-3-1-5-field-activity-fm-2025}.
  \item[\textbf{SLR components.}] Direct labor, indirect costs, and \ac{aor} recoupment~\autocite{edo-3-1-5-field-activity-fm-2025}.
  \item[\textbf{Document choice.}] A mission-funded activity using \ac{omn} in August for work starting in September should use a \ac{po} because the work completes after fund expiration~\autocite{edo-3-1-5-field-activity-fm-2025}.
\end{description}
%====================
% End of file
%====================
\IfSubfilesClassLoaded{
  \RenewDocumentCommand{\entryname}{}{\textbf{\color{Modern} Acronym}}
  \RenewDocumentCommand{\descriptionname}{}{\textbf{\color{Modern} Definition}}
  \printnoidxglossary[
    type=\acronymtype,
    title=Acronyms,
    style=long-booktabs
  ]}{}
\end{document}
